% ============================================================
%  LAMPIRAN 10 – PEMBAGIAN TUGAS DAN KONTRIBUSI ANGGOTA
%  Compile standalone:  pdflatex lampiran-10.tex
%  Compile as part of thesis: pdflatex main.tex
% ============================================================
\documentclass[../../thesis]{subfiles}
\usepackage{hyphenat}
\usepackage{iftex}

% ── Encoding, Language & Fonts ───────────────────────────────
\ifXeTeX
    \usepackage{polyglossia}
    \setmainlanguage{bahasai}
    \usepackage{fontspec}
    \setmainfont{TeX Gyre Termes}
    \setsansfont{TeX Gyre Heros}
    \setmonofont[Scale=0.88]{TeX Gyre Cursor}
\else
    \usepackage[utf8]{inputenc}
    \usepackage[T1]{fontenc}
    \usepackage[indonesian]{babel}
    \usepackage{mathptmx}
    \usepackage{helvet}
    \usepackage{courier}
\fi

% ── Page geometry (A4, UI thesis margins) ────────────────────
\usepackage[
  a4paper,
  top    = 2.54cm,
  bottom = 2.54cm,
  left   = 3.17cm,
  right  = 2.54cm
]{geometry}

% ── Line spacing & paragraph ─────────────────────────────────
\usepackage{setspace}
\onehalfspacing
\setlength{\parindent}{1.27cm}
\setlength{\parskip}{0pt}

% ── Microtype & justification ────────────────────────────────
\usepackage{microtype}
\sloppy

% ── Headings ─────────────────────────────────────────────────
\usepackage{titlesec}

\titleformat{\chapter}[block]
  {\centering\rmfamily\bfseries\large}{}
  {0pt}{\MakeUppercase}
\titlespacing*{\chapter}{0pt}{0pt}{1.5em}

\titleformat{\section}
  {\rmfamily\bfseries\normalsize}{\thesection}{1em}{}
\titlespacing*{\section}{0pt}{1.2em}{0.6em}

\titleformat{\subsection}
  {\rmfamily\bfseries\normalsize}{\thesubsection}{1em}{}
\titlespacing*{\subsection}{0pt}{1em}{0.5em}

% ── Tables ───────────────────────────────────────────────────
\usepackage{longtable, booktabs, array, multirow}
\usepackage{calc}
\usepackage{tabularx}

% ── Hyperlinks ───────────────────────────────────────────────
\usepackage[hidelinks]{hyperref}

% Bibliography setup is inherited from main.tex (biblatex + references.bib)
% when this file is compiled standalone via the subfiles class.

% ============================================================
\begin{document}

% When standalone: switch to appendix mode and set the counter so the
% appendix heading renders correctly. When compiled inside main.tex,
% main.tex issues \appendix and advances the counter before \subfile.
\ifSubfilesClassLoaded{}{\appendix\setcounter{chapter}{1}}

\lampiran{Pembagian Tugas dan Kontribusi Anggota}
\label{lamp:pembagian-tugas}

Penelitian ini dikerjakan secara berkelompok oleh tiga peneliti. Pada
tahap awal, ketiga anggota terlebih dahulu melakukan eksplorasi
menyeluruh secara mandiri terhadap keseluruhan alur \textit{pipeline},
mulai dari proses \textit{ingestion} dan ekstraksi teks dari berkas PDF
buku teks, eksperimen \textit{prompting} terhadap LLM, hingga konstruksi
dan visualisasi \textit{knowledge graph} pada basis data Neo4j. Hasil
eksplorasi masing-masing anggota kemudian didiskusikan bersama dosen
pembimbing. Melalui proses bimbingan tersebut, hasil konstruksi
\textit{knowledge graph} yang dikembangkan oleh Yhoga dipilih sebagai
dasar \textit{pipeline} pada tahap selanjutnya karena dinilai paling baik.

Setelah arah \textit{pipeline} ditetapkan, pengerjaan dibagi berdasarkan
tahapan \textit{pipeline} KGC. Yhoga berfokus pada tahap ekstraksi dan
konstruksi, meliputi eksperimen untuk menentukan proses ekstraksi teks
PDF yang paling baik, perancangan \textit{prompt} LLM yang menghasilkan
entitas dan relasi yang sesuai, hingga konstruksi dan visualisasi graf
pada Neo4j. Tegar menangani tahap evaluasi dan analisis relasi hasil
sistem. Untuk keperluan ini, ia mengembangkan aplikasi web KG Review App
sebagai instrumen pengumpulan penilaian pakar dan melakukan perekrutan
panel pakar penilai. Sesuai kualifikasi yang ditetapkan pada
Subbab~\ref{subsec:subjek-evaluasi}, pakar yang dilibatkan merupakan
mahasiswa aktif atau lulusan dengan latar belakang disiplin ilmu yang
linier dengan mata pelajaran terkait. Pada tahap ini Tegar melaksanakan
analisis validasi menggunakan metrik kesepakatan antarpenilai (Gwet's
AC1) serta metrik kinerja relasi (Fuzzy Precision, Fuzzy Recall, dan
F1-Score), menjalankan mekanisme \textit{LLM-as-a-judge} untuk
mengajudikasi ketidaksepakatan penilaian antarpakar dan menetapkan
\textit{gold standard} berdasarkan konteks buku sumber. Adapun Soros
mengerjakan tahap penyempurnaan (\textit{completion}/KGC), yaitu prediksi
relasi implisit lintas-buku antarkonsep menggunakan kemiripan
\textit{embedding} dan klasifikasi LLM. Ringkasan pembagian tugas
disajikan pada Tabel~\ref{tab:pembagian-tugas}.

\begin{longtable}{@{}
  >{\raggedright\arraybackslash}p{0.27\linewidth}
  >{\raggedright\arraybackslash}p{0.49\linewidth}
  >{\raggedright\arraybackslash}p{0.16\linewidth}@{}}
\caption{Pembagian Tugas Anggota Penelitian}
\label{tab:pembagian-tugas}\\
\toprule
\textbf{Tahap/Aktivitas} & \textbf{Deskripsi} & \textbf{Penanggung Jawab} \\
\midrule
\endfirsthead
\toprule
\textbf{Tahap/Aktivitas} & \textbf{Deskripsi} & \textbf{Penanggung Jawab} \\
\midrule
\endhead
\bottomrule
\endlastfoot
Eksplorasi awal \textit{pipeline}
  & Eksperimen mandiri menyeluruh: \textit{ingestion} dan ekstraksi teks
    PDF, \textit{prompting} LLM, konstruksi, serta visualisasi Neo4j;
    hasil terbaik dipilih melalui bimbingan
  & Yhoga, Tegar, Soros \\
\addlinespace
Ekstraksi teks \& \textit{knowledge extraction}
  & Eksperimen ekstraksi teks PDF terbaik dan perancangan \textit{prompt}
    LLM penghasil entitas serta relasi
  & Yhoga \\
\addlinespace
Konstruksi \& visualisasi KG
  & Pembentukan dan visualisasi \textit{knowledge graph} pada Neo4j
  & Yhoga \\
\addlinespace
Validasi \& analisis relasi
  & Pengembangan KG Review App sebagai instrumen penilaian panel pakar,
    perekrutan pakar, dan analisis validasi (Gwet's AC1; Fuzzy
    Precision, Fuzzy Recall, dan F1-Score)
  & Tegar \\
\addlinespace
\textit{LLM-as-a-judge} \& konsensus
  & Ajudikasi ketidaksepakatan penilaian antarpakar untuk menetapkan
    label konsensus final (\textit{gold standard}) berdasarkan konteks
    buku sumber
  & Tegar \\
\addlinespace
Penyempurnaan KG (\textit{completion}/KGC)
  & Prediksi relasi implisit lintas-buku menggunakan kemiripan
    \textit{embedding} dan klasifikasi LLM
  & Soros \\
\end{longtable}

\end{document}
